%% kapa: exam=1 round=26 year=2015 period=2 form=A subject=appraisal_law subject_ko=부동산관계법규 seq=11 number=11
%% source: Q-Net 감정평가사 기출문제 (cst00309) artlSeq=5138001 file=extracted/1008/5138001/제26회_감정평가사 1차_기출/제26회_감정평가사 1차_기출/2교시_감평_A형.pdf
%% conversion: kordoc 4.24.51 -> tools/gen1.py (auto); review=proofread-sonnet; flags=none
\begin{question}{11}
국토의 계획 및 이용에 관한 법령상 토지거래의 허가 등에 관한 설명으로 옳지 않은 것은?
\begin{choices}
  \item 허가구역의 지정은 허가구역의 지정을 공고한 날부터 7일 후에 그 효력이 발생한다.
  \item 허가구역에 있는 토지의 소유권을 이전하는 토지거래계약 당사자의 한쪽이「지방공기업법」에 의한 지방공사인 경우 그 기관의 장이 시장ㆍ군수 또는 구청장과 협의를 할 수 있고, 그 협의가 성립된 때에는 토지거래계약에 관한 허가를 받은 것으로 본다.
  \item 「공익사업을 위한 토지 등의 취득 및 보상에 관한 법률」에 따른 토지의 수용의 경우에는 토지거래계약에 관한 허가 규정이 적용되지 아니한다.
  \item 「민원사무처리에 관한 법률」에 따른 처리기간에 허가증의 발급 또는 불허가처분 사유의 통지가 없거나 선매협의 사실의 통지가 없는 경우에는 그 기간이 끝난 날의 다음날에 토지거래계약에 관한 허가가 있는 것으로 본다.
  \item 법령의 규정에 의하여 조세ㆍ부담금 등을 토지로 물납하는 경우에는 토지거래계약에 관한 허가 규정이 적용되지 아니한다.
\end{choices}
\end{question}
\answer{1}
